\documentclass{situationist-preprint}

\title{Finkish and Masked Dispositions}
\msid{MET-11-b}
\msfamily{MET-11}
\msdiscipline{Metaphysics}
\mstype{Companion manuscript}
\mskeywords{dispositions; finks; masks; antidotes; intrinsic finks; conditional analysis}

\begin{document}
\maketitle

\begin{abstract}
MET-11-a argues that dispositions are real and irreducible, partly on the ground that finks, masks and mimics separate every disposition from every conditional stated without dispositional vocabulary. We examine the cases in more detail. Finks and masks differ in where the interference acts: on the disposition itself, or on the process from stimulus to manifestation. Lewis's reformed analysis handles extrinsic finks by holding an intrinsic basis fixed; intrinsic finks and intrinsic masks defeat it, because the interference then lies within the basis. In every variant, the only principled description of the interfering factor is that it prevents the disposition from being manifested, which is the irreducibility MET-11-a claims.
\end{abstract}

\section{Two loci of interference}

A disposition, in the sense of MET-11-a (Definition~2.1), is a property $D$ whose ascription specifies a stimulus $S$ and a manifestation $M$, such that a bearer of $D$ is one that would $M$ if $S$, in the absence of interference. The simple conditional analysis drops the final clause. Lemma~3.2 of MET-11-a shows that no conditional without dispositional vocabulary can replace it. The counterexamples fall into two classes, distinguished by where the interference acts.

\paragraph{Interference with the disposition.} A \emph{fink} removes the disposition when the stimulus occurs, or, in the reverse case, confers it \citep{Martin1994}. The electro-fink attached to a live wire makes the wire dead when it is touched; the wire is live until then, and it would not conduct if touched. Here the conditional fails because the bearer no longer has the disposition by the time the manifestation would occur.

\paragraph{Interference with the process.} A \emph{mask} leaves the disposition in place and blocks the process from stimulus to manifestation \citep{Johnston1992}. The packed glass remains fragile when struck; the packing absorbs the blow. An \emph{antidote} is a mask that acts after the process has begun: the poison is ingested and remains lethal, and the antidote intercepts its action \citep{Bird1998}.

The two classes call for different repairs to the conditional analysis. Finks call for holding the bearer's condition fixed during the stimulus; masks call for excluding interfering factors in the circumstances.

\section{The reformed analysis and its limits}

\citet{Lewis1997} repaired the analysis against finks. On his reformed analysis, $x$ is disposed to $M$ when $S$ if and only if, for some intrinsic property $B$ that $x$ has, if $x$ were $S$ and retained $B$ for a sufficient time, $S$ and $x$'s having $B$ would jointly be an $x$-complete cause of $x$'s $M$-ing. Since the fink is external and acts by removing $B$, holding $B$ fixed neutralizes it.

Two variants defeat the repair.

\paragraph{Intrinsic finks.} Suppose the device that removes the disposition is part of the bearer. A wire has, within its own insulation, a mechanism that detects contact and severs its conducting core; or an organism's immune system responds to an infection by eliminating the very susceptibility the infection exploits. Then the basis $B$ that the analysis holds fixed either includes the mechanism, in which case holding $B$ fixed includes holding fixed the mechanism that removes the disposition and the counterfactual still fails, or excludes it, in which case $B$ is not the bearer's total relevant intrinsic condition and the choice of $B$ has no principled ground other than that it is the basis of the disposition \citep{Clarke2008}.

\paragraph{Intrinsic masks.} Suppose the mask is part of the bearer: a glass with a protective inner lattice that absorbs blows of the stimulus magnitude without affecting the glass's molecular structure. The glass retains its basis, is struck, and does not break. Holding $B$ fixed does not help, since $B$ is not removed; the process is blocked inside the bearer. If one says that this glass is not fragile, one must say what distinguishes it from a glass packed in an external box, which is fragile; and the only difference is the location of the mask, which is irrelevant to whether the glass is constituted so as to break when struck.

\begin{proposition}\label{p:interference}
In every variant, the condition that distinguishes bearers of the disposition from non-bearers, among things that would or would not $M$ if $S$, is the presence or absence of a factor that prevents the disposition from being manifested.
\end{proposition}

\begin{proof}
Extrinsic finks and masks are distinguished from genuine absences of the disposition by the presence of an external factor that removes the disposition or blocks its manifestation. Intrinsic finks and masks are distinguished in the same way, by an internal factor that does the same. Mimics are distinguished by the absence of the disposition despite the manifestation. In each case, the factor is identified by its role, as what removes, blocks or mimics the disposition, and not by any intrinsic or extrinsic feature it has independently of that role: packing, inner lattices, immune responses and sorcerers have nothing else in common.
\end{proof}

Proposition~\ref{p:interference} is the detail behind Lemma~3.2 of MET-11-a. Every conditional that survives the cases contains the notion of interference with the disposition, and interference is not specifiable except by reference to the disposition interfered with.

\section{Objections}

\begin{objection}[Habitual analysis]
\citet{Fara2005} analyzes dispositions by habituals: $x$ is fragile if $x$ breaks when struck, in the sense of a generic or habitual sentence, which tolerates exceptions. Masks and finks are exceptions of the kind habituals tolerate.
\end{objection}

\begin{reply}
The habitual analysis is not a reduction in the sense of MET-11-a (Definition~2.3). Habitual sentences are true of things that would display the behavior in normal conditions, and what makes conditions normal for a particular bearer is that nothing prevents its dispositions from being manifested. The analysis therefore reports the irreducibility rather than removing it; and it is consistent with the present thesis if taken as an account of the logical form of disposition ascriptions rather than as a reduction.
\end{reply}

\begin{objection}[Specify the stimulus]
\citet{Choi2006} argues that the simple analysis survives if the stimulus is specified finely enough: striking the glass when it is unpacked, touching the wire when no fink is attached.
\end{objection}

\begin{reply}
Specifying the stimulus case by case works only for the finite list of interferers already considered. For the intrinsic variants, the specification would have to exclude a condition of the bearer itself, and the only general way to say which conditions to exclude is to say that they are those that prevent the disposition from being manifested. That specification is not free of dispositional vocabulary.
\end{reply}

\section{Conclusion}

Finks interfere with the disposition, masks with its manifestation, and both may lie inside the bearer as well as outside it. The reformed analysis meets the external fink and fails at the internal fink and the internal mask. What unites the cases is a role, prevention of manifestation, which can be specified only by reference to the disposition prevented. Dispositions are therefore not reducible to conditionals, and, as MET-12-a argues, they belong to the mosaic that the best system of laws summarizes rather than being grounded in those laws.

\begin{thebibliography}{9}\small
\bibitem[Bird(1998)]{Bird1998} Bird, A. (1998). Dispositions and antidotes. \emph{Philosophical Quarterly}, 48(191), 227--234.
\bibitem[Choi(2006)]{Choi2006} Choi, S. (2006). The simple vs. reformed conditional analysis of dispositions. \emph{Synthese}, 148(2), 369--379.
\bibitem[Clarke(2008)]{Clarke2008} Clarke, R. (2008). Intrinsic finks. \emph{Philosophical Quarterly}, 58(232), 512--518.
\bibitem[Fara(2005)]{Fara2005} Fara, M. (2005). Dispositions and habituals. \emph{No\^us}, 39(1), 43--82.
\bibitem[Johnston(1992)]{Johnston1992} Johnston, M. (1992). How to speak of the colors. \emph{Philosophical Studies}, 68(3), 221--263.
\bibitem[Lewis(1997)]{Lewis1997} Lewis, D. (1997). Finkish dispositions. \emph{Philosophical Quarterly}, 47(187), 143--158.
\bibitem[Martin(1994)]{Martin1994} Martin, C.~B. (1994). Dispositions and conditionals. \emph{Philosophical Quarterly}, 44(174), 1--8.
\end{thebibliography}

\end{document}
